주어지는 난해한 프로그래밍 언어 「Brainfxxk」의 소스코드의 실행 결과를 작성하여라. 소스코드의 길이는 $10,000$ 를 넘지 않는다.  

\\

Brainfxxk는 아래의 규칙에 따라 실행된다:

\begin{enumerate}
\item 
Brainfxxk 인터프리터는 32768개 크기의 \texttt{char} 배열과 이 배열의 인덱스를 가리키는 포인터를 메모리 공간으로 가진다. 배열의 모든 요소와 포인터는 \texttt{0}으로 초기화된다.

\begin{lstlisting}
// define
char mem[32768];
int ptr;

// init
memset(mem, 0, sizeof(mem));
ptr = 0;
\end{lstlisting}

\item 
인터프리터는 소스코드의 각 문자를 왼쪽 끝에서부터 오른쪽 방향으로 읽는다. 오른쪽 끝 문자까지 읽었다면, 다음 줄의 왼쪽 끝에서부터 같은 행동을 반복한다.


\item 
2 과정에서 읽어들인 개별 문자는 아래와 같이 처리한다.

\begin{itemize}
\item 
\texttt{>} : 포인터 1 증가\\
32768 이상의 메모리 주소는 존재하지 않는다. 메모리 주소 32768은 오버플로우되어 메모리 주소 0이다.


\item 
\texttt{<} : 포인터 1 감소\\
음의 메모리 주소는 존재하지 않는다. 메모리 주소 -1은 언더플로우되어 메모리 주소 32767이다.


\item 
\texttt{+} : 포인터가 가리키는 요소의 값을 1 증가


\item 
\texttt{-} : 포인터가 가리키는 요소의 값을 1 감소


\item 
\texttt{.} : 포인터가 가리키는 요소의 값을 아스키 코드 문자로 출력


\item 
\texttt{,} : 표준입력으로부터 값을 입력받아, 그 값의 아스키 코드를 포인터가 가리키는 요소에 덮어씀


\item 
\texttt{[} - 포인터가 가리키는 바이트의 값이 0이라면 짝이 되는 \texttt{]}로 이동


\item 
\texttt{]} - 포인터가 가리키는 바이트의 값이 0이 아니면 짝이 되는 \texttt{[}로 이동


\item 
\texttt{\0} : EOF(End of File); 프로그램을 종료한다.


\item 
실제 Brainfxxk과 달리, 이 문제에서는 \texttt{,} 문자는 무시한다.


\item 
위 정의를 제외한 모든 문자는 무시하고 계속 진행한다.


\end{itemize}

\end{enumerate}

인터프리터는 EOF(\texttt{\0}; 아스키코드 0)를 입력받기 전까지는 계속해서 입력을 받아와야 한다. 만약 로컬에서 답안을 실행하는 도중, 파이프라인¹을 사용하지 않고 직접 표준입력에 값을 입력하고자 한다면, EOF 문자를 직접 입력하여야 한다.

\begin{itemize}
\item 
윈도우: \texttt{Ctrl} + \texttt{Z}를 입력하고 \texttt{Enter} 한다.


\item 
리눅스/mac: \texttt{Ctrl} + \texttt{D}를 입력하고 \texttt{Enter} 한다.


\end{itemize}

이 과정에서 표준 입력에 입력으로 나타나는 \texttt{^Z}와 \texttt{^D}는 실제 \texttt{^Z}, \texttt{^D}라는 문자열이 아닌, \texttt{Ctrl}과 \texttt{Z} 혹은 \texttt{D}를 조합해 입력한 문자라는 의미이다.


¹ \texttt{cat <input> | <compiled-program>}

\\

아래는 몇 가지 Brainfxxk 예제이다.

\\

예제 입력 1

\begin{lstlisting}
+++++++++++++++++++++++++++++++++.
\end{lstlisting}

예제 출력 1

\begin{lstlisting}
!
\end{lstlisting}

예제 1에서 \texttt{+}가 33개 등장하여 0번 메모리 공간의 값이 33이 되었다. \texttt{.}에 의해 아스키코드 33이 나타내는 문자, \texttt{!}가 출력되었다.

\\

예제 입력 2

\begin{lstlisting}
+++++[>++++++++++<-]>.
\end{lstlisting}

예제 출력 2

\begin{lstlisting}
2
\end{lstlisting}
\begin{enumerate}
\item 
앞의 \texttt{+++++}로 0번 메모리 공간의 값이 5가 되었다. \texttt{[}를 처리할 때, 포인터가 가리키는 0번 메모리 공간의 값이 0이 아니므로 \texttt{]}로 건너뛰지 않고 계속 진행한다.


\item 
\texttt{[} 다음의 \texttt{>}에 의해 포인터가 1 증가하여, 이후부터는 1번 메모리 공간을 사용한다.


\item 
\texttt{>} 다음에 \texttt{+}가 10개 등장하여, 1번 메모리의 값은 10이 되었다.


\item 
\texttt{+} 다음의 \texttt{<}에 의해 포인터가 1 감소하여, 이후부터는 0번 메모리 공간을 사용한다.


\item 
\texttt{<} 다음의 \texttt{-}에 의해 0번 메모리 공간의 값이 1 감소한 4가 되었다.


\item 
\texttt{]}를 처리할 때, 포인터가 가리키는 0번 메모리 공간의 값이 0이 아니므로, 계속 진행하지 않고 \texttt{[}로 되돌아간다. (\texttt{mem = { 4, 10, ... };})


\item 
2 ~ 6 과정을 4회 더 반복하여 \texttt{mem = { 0, 50, ... };}이 된다.


\item 
\texttt{]}를 처리할 때, 포인터가 가리키는 0번 메모리 공간의 값이 0이므로, 더 이상 되돌아가지 않고 다음 문자를 처리한다.


\item 
\texttt{>}에 의해 포인터가 1 증가하여, 1번 메모리 공간을 사용한다.


\item 
\texttt{.}에 의해 아스키코드 50이 나타내는 문자, \texttt{2}가 출력된다.


\end{enumerate}
